\documentclass[11pt]{article}
\usepackage{series}
\usepackage{mathrsfs}

\title{Closure Geometry and a Regime-Local\\
Ten-Dimensional Action Ansatz\\
\large Corrected fourth edition}
\author{Peter Nero}
\date{July 2026}

\begin{document}
\maketitle

\begin{abstract}
We formulate a ten-dimensional effective action compatible with closure-strain
and q79 carrier data.  The metric, bundle $M_{10}\to Y_4$, compact fiber,
field representations, and derivative expansion are declared realization
inputs.  The action is an ansatz on a specified regime, not the most general
action and not a derivation of gravity from projection.  We give the conditions
under which a closure Hessian contributes to canonically normalized pole
masses, a rank-one alignment projector yields one Higgs doublet, an internal
operator has discrete spectrum, and a four-dimensional mode truncation is
consistent.  Curvature is defined through explicit connections rather than
inferred from nonuniform strain.  The selected q79 Fu--Yau branch is the
compactification candidate, while Lens--Nil remains auxiliary.  Existing
finite-matrix and Standard-Model profile calculations are incorporated at
their declared embedded-renormalized-SM tier; their profile inputs and imported
BRST quantization are not reclassified as no-knob consequences of this action.
\end{abstract}

\section*{Revision note for this edition}
\addcontentsline{toc}{section}{Revision note for this edition}
\begin{description}
\item[Supersedes.] \emph{Closure Geometry and Unified Dynamics: A
Ten-Dimensional Action for Mass, Scalar Relaxation, Quantization, and
Curvature}, version~3.
\item[Reason.] The metric and Einstein--Hilbert term were imported while the
action was described as derived and most general; closure cost, Hessian
positivity, nil language, and nonuniform strain were also promoted directly to
mass, one Higgs, quantization, and curvature.
\item[Resolution.] Version~4 declares a regime-local EFT ansatz, lists omitted
operators, and supplies separate pole-mass, alignment-projector,
compact-resolvent, curvature, and consistent-truncation gates.
\item[Retained result.] The action remains a useful conditional synthesis and
can host the currently certified finite-SM profile branch.
\item[Remaining boundary.] A same-source q79 action, normalized zero modes,
gauge-fixed Hessians, reduction error, and strict value selection remain to be
constructed.
\end{description}

\tableofcontents

\section{Status and regime}

Fix an energy interval $E\ll\Lambda_{10}$ and an admissible field neighborhood
$\mathcal U$ on which a local derivative expansion is valid.  The word
``minimal'' below means that only a displayed subset of operators is retained
for a stated calculation.  It does not mean uniqueness under all symmetries.

The action requires as input:
\begin{enumerate}
\item a ten-dimensional smooth bundle
$\pi:M_{10}\to Y_4$ with compact six-dimensional fiber $X_x$;
\item a slab-local Lorentzian metric $g_{10}$ with one time direction and a
globally hyperbolic four-dimensional physical base;
\item gauge and spin/$\operatorname{Spin}^c$ bundles with compatible
connections;
\item the selected q79 Fu--Yau geometry and HYM data, or another explicitly
declared compactification;
\item a field content and symmetry group; and
\item a power-counting rule and cutoff controlling omitted operators.
\end{enumerate}

The component identity $1+3\times3=4+6$ motivates a local comparison carrier
but supplies none of these global inputs.

\section{Fields and typed geometry}

Let $Q_{\rm WW}\in\Gamma(\operatorname{Hom}(TP,TI))$ be the rank-three
comparison field and let
\[
 S=\frac12\log(Q_{\rm WW}^TQ_{\rm WW})
\]
be its local strain.  A selected flag gives the strain projectors
$P_{\rm sc},P_{\rm sh},P_{\rm nil}$ of ranks $1,2,3$.

For the q79 degree-three cover, the selected internal carrier is
\[
 \mathcal H_{\rm q79}
 =L_{\rm shared}\otimes
 (\mathcal O\oplus\mathcal A_0\oplus\mathcal A),
 \qquad \operatorname{rank}=1+2+3.
\]
An action using $S$ as the source of q79 fields must include or cite a global
intertwiner $\mathfrak I_{\rm WW\to q79}$ preserving metrics, connections, and
the selected operators.  This paper does not infer it from rank equality.

Let $\phi^A$ denote real scalar coordinates on the retained field manifold,
$A_M^a$ a gauge connection, $\Psi$ a spinor in a declared representation, and
$H$ a selected complex scalar module when present.

\section{Regime-local action ansatz}

One two-derivative ansatz is
\begin{align}
S_{10}=\int_{M_{10}}\!\sqrt{|g_{10}|}\,\dd^{10}x\,
\Big[&\frac{1}{2\kappa_{10}^2}R_{10}
-\frac14 k_{ab}(\phi)F^a_{MN}F^{b\,MN}
-\frac12G_{AB}(\phi)D_M\phi^A D^M\phi^B\nonumber\\
&-V(\phi)
+i\overline\Psi\Gamma^M D_M\Psi
-\overline\Psi\,\mathcal M(\phi)\Psi
+\mathcal L_{H,B,\Phi}
\Big]
+S_{\rm gf}+S_{\rm gh}+S_{\rm bdy}.
\label{eq:action}
\end{align}
Here $k_{ab}$ and $G_{AB}$ must be positive on physical directions,
$\mathcal L_{H,B,\Phi}$ records the selected flux/torsion/dilaton sector, and
the gauge-fixing, ghost, and boundary terms are part of the definition of a
quantized perturbative calculation.

Equation~\eqref{eq:action} imports the Einstein--Hilbert term and a Lorentzian
metric.  It therefore realizes gravity; it does not derive gravity from strain
or projection.

\subsection{Omitted operators}

Unless forbidden by a stated symmetry, the effective action can also contain
\[
 R^2,\quad R_{MN}R^{MN},\quad R_{MNPQ}R^{MNPQ},\quad
 RF^2,\quad F^3,\quad F^4,\quad (D\phi)^4,\quad
 R(D\phi)^2,\quad \phi^n,\quad
 \overline\Psi\Gamma\Psi D\phi,
\]
torsion and Chern--Simons terms, higher fermion operators, and higher
derivatives.  A truncation must bound their contribution by powers of
$E/\Lambda_{10}$ and the relevant curvature or field amplitudes.

Higher-curvature or higher-time-derivative terms can introduce extra degrees
of freedom or ghosts if treated nonperturbatively.  In an EFT treatment they
are perturbative operators below the cutoff; any claimed fundamental
completion requires a separate constraint/propagator analysis.

\section{Closure Hessian and physical masses}

Let $\varphi=0$ be a stationary background for retained real fields and write
the quadratic four-dimensional action after integrating the internal fiber as
\[
 S_4^{(2)}
 =\frac12\int_{Y_4}\sqrt{|g_4|}\,dd^4x
 \left(
 Z_{AB}\,\partial_\mu\varphi^A\partial^\mu\varphi^B
 -M^2_{AB}\varphi^A\varphi^B
 \right).
\]

\begin{theorem}[Pole-mass promotion]
Assume $Z$ is positive definite, the quadratic operator is self-adjoint on its
declared domain, and interactions admit the stated perturbative pole
definition.  At tree level the squared masses are the eigenvalues of
\[
 Z^{-1/2}M^2Z^{-1/2}.
\]
A closure Hessian $H_{\mathcal J}$ contributes to physical mass only after a
source theorem identifies $M^2=C^*H_{\mathcal J}C$ for a normalized field map
$C$.
\end{theorem}

\begin{proof}
The field redefinition $\widehat\varphi=Z^{1/2}\varphi$ canonically normalizes
the kinetic term.  The inverse propagator is then
$p^2I-Z^{-1/2}M^2Z^{-1/2}$, whose zeros give the tree-level poles.  The final
statement follows because an unnormalized dimensionless cost is not the
quadratic coefficient in the physical action until $C$ and units are fixed.
\end{proof}

Loop-corrected pole masses require the renormalized self-energy and a declared
mass scheme.  Eigenvalues of an internal Hessian are not automatically pole
masses.

\section{Scalar alignment and the Higgs gate}

A positive Hessian on the six-dimensional strain sector proves local
stability, not one scalar.  Suppose instead that the selected finite
fluctuation space $V_{\rm scal}$ carries the required gauge representation and
there is a source-selected orthogonal projector
\[
 P_H:V_{\rm scal}\to V_H,
 \qquad \operatorname{rank}_{\mathbb R}V_H=4,
\]
onto one complex weak doublet.

\begin{proposition}[Conditional one-doublet reduction]
If $P_H$ commutes with the quadratic gauge-fixed operator, the orthogonal
complement is massive above the truncation gap or consistently removed, and
the potential restricted to $V_H$ has the Standard-Model symmetry-breaking
form, then the retained scalar sector contains one Higgs doublet at that
scale.
\end{proposition}

The selected finite-algebra packet executes such a rank-four alignment
projector on a raw rank-twelve real fluctuation space and removes eight real
scalar directions at the declared profile tier.  That is the relevant
selection result.  It must not be replaced by ``the Hessian has a unique
radial direction,'' which is false without the projector and representation
data.

\section{Discrete spectra and quantization}

Let $D_X$ be a self-adjoint internal Dirac-type or Laplace-type operator on the
compact fiber with elliptic boundary conditions when a boundary is present.

\begin{theorem}[Internal spectral discreteness]
If $(D_X-i)^{-1}$ is compact, then the spectrum of $D_X$ is discrete with
finite-multiplicity eigenvalues accumulating only at infinity.
\end{theorem}

This theorem justifies a Kaluza--Klein or finite spectral expansion.  A nil
group, nil boundary, or divergent cost does not by itself imply compact
resolvent or isolated minima.  Moreover, spectral discreteness is not a
derivation of quantum probabilities.  Quantization still requires a state
space, observable algebra, dynamics, constraints, and probability rule.

\section{Curvature and integrability}

Let $P_H$ now denote a projector defining a horizontal distribution in a
declared bundle, and let $\nabla$ be a connection.  Its curvature is
\[
 F_\nabla=\nabla^2,
\]
or locally $F=\dd A+A\wedge A$.  Frobenius failure of a distribution is tested
by the vertical part of $[X,Y]$ for horizontal vector fields $X,Y$.

The condition $\nabla S\ne0$ says only that strain is nonparallel.  It does not
by itself imply nonintegrability, Riemann curvature, or the Einstein equations.
Any curvature--strain coupling in~\eqref{eq:action} must be written with an
explicit connection and varied to obtain its field equations.

\section{q79 Fu--Yau specialization}

The current selected global compactification candidate is the q79 Fu--Yau
branch with its declared complex, flux, spectral-cover, and HYM data.  Known
Hull--Strominger/Fu--Yau mathematics supplies a legitimate class of heterotic
backgrounds once anomaly cancellation and the relevant bundle conditions are
satisfied.

The auxiliary geometry $L(3,1)\times\mathrm{Nil}_3$ can support model
calculations but is not the same manifold: its cohomology already differs from
the q79 branch.  The action must choose one background and integrate over that
background.  Literal $S^1\times\mathrm{Lens}\times\mathrm{Nil}$ and literal
nesting are not used to define $X_6$.

The shared circle is encoded by $L_{\rm shared}$ and its connection.  It is
counted once and is not a time coordinate.

\section{Four-dimensional reduction}

Let $\{\chi_n(y)\}$ be normalized internal modes and expand a field as
\[
 \Phi(x,y)=\sum_{n\in I_{\rm keep}}\phi_n(x)\chi_n(y)
 +\sum_{r\in I_{\rm disc}}\eta_r(x)\chi_r(y).
\]
Setting all $\eta_r$ to zero is a consistent truncation only if their exact
Euler--Lagrange equations vanish on the retained ansatz:
\[
 \left.\frac{\delta S_{10}}{\delta\eta_r}\right|_{\eta=0}=0
 \quad\text{for every }r\in I_{\rm disc}.
\]

\begin{theorem}[Conditional coherent reduction]
Assume the retained mode space is invariant under all nonlinear terms to the
order considered, the discarded equations vanish on the retained ansatz, the
discarded linear operator has gap $\lambda_{\rm gap}>0$, and the nonlinear
source into discarded modes obeys $\|J_{\rm disc}\|\le\epsilon$.  Then the
leading eliminated field obeys
\[
 \|\eta\|\le
 \|L_{\rm disc}^{-1}\|\,\epsilon
 \le \frac{C\epsilon}{\lambda_{\rm gap}},
\]
and substituting it gives a controlled local four-dimensional effective action
with the corresponding Schur--Feshbach correction.
\end{theorem}

\begin{proof}
Solve the discarded equation by the inverse on the gapped complement and use
the stated resolvent bound.  The effective correction follows by substitution
or the Schur complement.  Exact consistency is the special case
$J_{\rm disc}=0$.
\end{proof}

A spectral gap alone does not prove nonlinear invariance or exact truncation.

\section{Gauge fixing, anomalies, and observables}

A perturbative gauge theory built from~\eqref{eq:action} must include a common
gauge-fixing, ghost, zero-mode, regulator, scale, and renormalization-scheme
policy.  BRST nilpotency and anomaly cancellation are equations to verify.
Physical predictions require a functor from action parameters to renormalized
local or scattering observables.

The current embedded-renormalized-SM closure imports standard
BRST/Faddeev--Popov quantization at its declared profile tier.  This is a valid
reconstruction standard, but it is not a derivation of the BRST/path-integral
or Born-record rules from MTT.

\section{Relation to the current numerical closure}

At the adopted one-shared-physical-primitive/profile standard, the calculation
repositories lock the finite $27\times27$ matrix, the physical finite Dirac
operator at profile tier, charged Yukawa magnitudes, CKM rows, electroweak and
threshold rows, multi-loop precision transport, and the renormalized
observable functor.  The final audit reports embedded renormalized-SM
equivalence at that stated standard.

This action paper neither demotes nor strengthens those certificates.  It
records what would be needed to derive the same data from one ten-dimensional
source action:
\begin{enumerate}
\item the world-in-world/q79 bundle-and-connection intertwiner;
\item the selected normalized internal zero modes and overlap kernels;
\item one common gauge-fixed Hessian and transport policy;
\item the consistent-truncation or error theorem; and
\item independent source selection if strict no-knob closure is claimed.
\end{enumerate}

\section{Scoped action theorem}

\begin{theorem}[Regime-local action statement]
Given the geometric, representation, metric, gauge-fixing, and EFT inputs of
this paper, action~\eqref{eq:action} defines a local covariant realization on
the declared domain.  Under compact-resolvent, alignment-projector,
pole-normalization, and consistent-reduction hypotheses, it yields a discrete
internal mode expansion, a selected one-doublet scalar sector, canonically
defined tree-level masses, and a controlled four-dimensional effective action.
\end{theorem}

The theorem is conditional.  It does not establish uniqueness of the action,
derive the metric or dimension, select the q79 vacuum, or prove the full
Standard Model without the listed source and observable maps.

\section{Conclusion}

The corrected action is a useful synthesis ansatz, not a universal derivation.
Its value is that every physical promotion now has a recognizable mathematical
gate: global geometry, connection, compact resolvent, scalar projector,
canonical pole normalization, gauge consistency, and controlled reduction.
The next decisive construction is the same-source q79 intertwiner followed by
evaluation of the action's normalized internal rows.

\begin{thebibliography}{9}
\bibitem{WeinbergEFT}
S.~Weinberg, \emph{The Quantum Theory of Fields, Vol. II}, Cambridge
University Press, 1996.

\bibitem{FuYau}
J.-X.~Fu and S.-T.~Yau,
\emph{The theory of superstring with flux on non-Kahler manifolds and the
complex Monge--Ampere equation}, J. Differential Geom. 78 (2008).

\bibitem{MTTSM}
P.~Nero, \emph{MTT Current True SM Closure Consolidated Ledger}, internal
theorem and verification packet, 2026.
\end{thebibliography}

\end{document}
